\section{Discussion}
\label{sec:discussion}

\paragraph{Incidence and intersection.} An incidence lookup in HStore reads one entry of the reverse index. For most atoms the whole incidence set is stored inline in that entry, and for hubs it is the root of a small tree. The visitor API hands each hyperedge id to the caller without allocating, which is why \wl{read.incidence} allocates about one byte per lookup. HyperGraphDB opens a JE cursor on a duplicate set, decodes 16-byte handles, and builds a result object per lookup. With a 64\,MiB JE cache, part of each lookup also misses the cache and reads from the log files (\cref{app:full}). Co-membership adds a second effect. HStore intersects two sorted incidence trees and skips every subtree whose key range cannot match, so the cost follows the smaller set. HyperGraphDB plans \texttt{hg.and} as a general conjunctive query over two result sets. The two-hop workload combines both lookups, so it inherits both advantages.

\paragraph{Member enumeration.} HyperGraphDB keeps a link's targets as a Java array in its atom cache, so enumerating the members of a cached link is an array loop with almost no allocation. HStore reads a hyperedge's members from its record in the catalog tree: from the leaf stored inline in the record for a small edge, or from the edge's own membership tree. What decides the result at scale 4 is how much heap a cached catalog leaf takes. A leaf holds about 80 records, each with its inline member list. Decoded in full, it takes roughly 11 times the bytes of its page image, and at scale 4 only about half of the catalog then fits in a 64\,MiB node cache. An earlier HStore version (commit \texttt{459b8be}) decoded leaves in full, and its member enumeration at scale 4 ran at 0.13 times HyperGraphDB's rate, with most of the time spent decoding records that the lookup did not need. The version measured here frames each record by its length and decodes a record only when it is read, so a cached leaf costs its page bytes plus the records actually used. The catalog leaves that the lookups reach then stay in the cache, and member lookups read nothing from the files.

\paragraph{Tail latency.} HStore readers pin an immutable generation and never wait for a writer, a checkpoint or a compaction. Writers wait only for the commit lock, which a checkpoint takes just to swap stored references into the head. HyperGraphDB reads run inside JE transactions, which take record locks, and JE's cleaner and checkpointer compete for the same latches and log. In 16 of the 18 latency measurements (six workloads in three configurations), the gap between the two systems is wider at p99 and p99.9 than at the median. The two exceptions are the \wl{mixed} writer under synchronous commits and at scale 4. There the gap is widest at p99 and smaller at p99.9 than at the median, and since the writer makes 1,000 commits per run, its p99.9 is its slowest commit.

\paragraph{Giant hyperedges.} HStore bulk-builds a large membership tree bottom-up in one pass, and a membership probe is one descent of that tree. HyperGraphDB stores the giant link as one record holding 50,000 handles (200,000 at scale 4). A membership probe goes through the node's incidence set instead.

\paragraph{Footprint and reclamation.} Several layout choices keep HStore's files small: 8-byte ids stored as varint gaps instead of 16-byte handles, node images packed into 64-byte units, and small member trees and incidence sets stored inline in their parent record. HStore's log is also truncated at each checkpoint. A JE log keeps every superseded version until the cleaner copies the live entries out of a file and deletes it, and at scale 4 cleaning dominated HyperGraphDB's bytes written. HStore's weak point is the period right after a large delete. Pages replaced by the deletes stay on disk while any of the 64 retained generations can reach them, so the files grow first and shrink only after the history moves on, here at the restart that follows. An application that deletes heavily and needs its disk space back quickly would set a smaller \texttt{history\_limit}.

\paragraph{Opening and recovery.} After a clean shutdown, HStore reads one catalog image and finds nothing in the log to replay. HyperGraphDB's open takes longer and allocates far more heap: \M{allocAhgdbreopen} per reopen at scale 1, against \M{allocAhstorereopen} for HStore. After a crash, HStore replays the logged key changes of every commit since its last checkpoint, and JE replays its log from its last checkpoint. HStore's replay cost grows with the number of commits since the last checkpoint, which the write buffer and the checkpoint thresholds bound. At scale 1 HyperGraphDB opens the crashed store faster, \M{gArecover} against \M{hArecover}. HStore's time there is close to its clean reopen of the full dataset (\M{hAreopen}), and it allocates \M{allocAhstorerecover} while opening, against \M{allocAhgdbrecover} for HyperGraphDB. At scale 4 the order reverses: \M{hCrecover} for HStore against \M{gCrecover}.

\paragraph{What HStore spends.} HStore retains more heap after ingest than HyperGraphDB (\cref{tab:memory}). The retained heap includes nodes not yet written by a checkpoint, up to the 64\,MiB write buffer, and the node cache, whose 64\,MiB budget counts estimated rather than measured node sizes. At scale 1 it also allocates more per inserted node and hyperedge, opens a crashed store more slowly, and with asynchronous commits its median single-update commit is slightly slower than HyperGraphDB's. Its slowest single commit at scale 4 was slower than HyperGraphDB's, even though its p99 and p99.9 were lower.
